\clearpage
\chapter*{第232次报告勘误}
\markboth{第232次报告勘误}{第232次报告勘误}
\addcontentsline{toc}{chapter}{第232次报告勘误}
\sid{e232}
\src{303}

\noindent\textbf{7. TDTE V：Picard 概形。存在性定理。}

\begin{center}
[布尔巴基讨论班，第14卷，1961/62年度，第232号，19页。]
\end{center}

\noindent\textbf{第7页，评注3.3。} --- 这里提出的问题已由 Mumford 作出
肯定解答（见下一篇报告的评注）。

\noindent\textbf{第13页，评注5.2。} --- 正如我在下一篇报告开头所指出的，
这里提出的存在性猜想是错误的；不过 Mumford 用他的方法证明了一个稍弱的
定理。

\noindent\textbf{第14页，第7行。} --- 在“完备”之后加上“具有代数闭剩余
域”。刚才提到的 Mumford 反例还表明，这一限制不可缺少。

\noindent\textbf{第17页，评注6.6。} --- 正如我们在下一篇报告开头所指出的，
这里提出的问题刚刚由 Murre 作出肯定解答。

\noindent\textbf{第18页，评注6.7。} --- 这里提出的问题有肯定解答（参见下一
篇报告评注的最后一段）。

\noindent\textbf{第18页，评注6.8，第4行。} --- 把“正则”改为“在 $k$ 上
光滑”。

\par
